\documentclass[../../main.tex]{subfiles}% 注意这里的文件路径不能用 ./main.tex ,否则用latexmk编译子文件会报错
\graphicspath{{\subfix{./image/}}{\subfix{../../image/}}} % 指定图片引用路径,后续可以直接使用图片文件名
% 如果图片存放在上级目录,则使用 ../../image/ 作为备用路径

% 例如:
% \begin{figure}[H]
% \centering
% \includegraphics[width=0.3\textwidth]{图.png}
% \caption{}
% \label{figure:图}
% \end{figure}
% 注意:上述\label{}一定要放在\caption{}之后,否则引用图片序号会只会显示??.

\begin{document}

\section{习题}

\begin{proposition}
假设 $f:\mathbb{R}^n\to\mathbb{R}$ 是无穷次可微函数，证明：
\begin{align*}
f(x) = \sum_{|\alpha|\leqslant k} \frac{1}{\alpha!} D^\alpha f(0)x^\alpha + O(|x|^{k+1}),
\end{align*}
其中 $k$ 可为任何正整数且 $\alpha=(\alpha_1,\alpha_2,\cdots,\alpha_n)$ 表示多重指标，$|\alpha|=\alpha_1+\alpha_2+\cdots+\alpha_n$，$\alpha!=\alpha_1!\alpha_2!\cdots\alpha_n!$，$x^\alpha=x_1^{\alpha_1}x_2^{\alpha_2}\cdots x_n^{\alpha_n}$。
\end{proposition}

\begin{example}
求二阶常微分方程 $y''=y^3-y$ 的一个特解。
\end{example}

\begin{proposition}
证明 Leibniz 公式
\begin{align*}
D^\alpha(uv) = \sum_{\beta\leqslant \alpha} \binom{\alpha}{\beta} D^\beta u\,D^{\alpha-\beta}v,
\end{align*}
其中 $\alpha=(\alpha_1,\alpha_2,\cdots,\alpha_n)$，$\beta=(\beta_1,\beta_2,\cdots,\beta_n)$ 表示多重指标，
\begin{align*}
\binom{\alpha}{\beta} = \frac{\alpha!}{\beta!(\alpha-\beta)!},
\end{align*}
且 $\beta\leqslant\alpha$ 表示 $\beta_i\leqslant\alpha_i$ $(i=1,2,\cdots,n)$。
\end{proposition}

\begin{example}
讨论二阶偏微分方程
\begin{align*}
A u_{xx} + 2B u_{xy} + C u_{yy} = f(x,y,u_x,u_y)
\end{align*}
的分类，其中 $u=u(x,y)$；$A,B,C$ 是参数。
\end{example}

\begin{example}
将下列方程化为标准型：
\begin{enumerate}[(1)]
\item $u_{xx} + 2u_{xy} + 2u_{yy} = 0$；
\item $u_{xx} + 2u_{xy} + u_{yy} = 0$；
\item $u_{xx} + 4u_{xy} + u_{yy} = 0$.
\end{enumerate}
\end{example}

\begin{example}
将下列方程化为标准型：
\begin{enumerate}[(1)]
\item $\sum_{i=1}^n u_{x_i x_i} + \sum_{1\leqslant i < j \leqslant n} u_{x_i x_j} = 0$；
\item $\sum_{1\leqslant j < i \leqslant n} u_{x_i x_j} = 0$.
\end{enumerate}
\end{example}

\end{document}